\documentclass[11pt,a4paper]{article}
\usepackage[utf8]{inputenc}
\usepackage[T1]{fontenc}
\usepackage{lmodern}
\usepackage{amsmath}
\usepackage[margin=2cm]{geometry}
\usepackage{enumitem}
\DeclareUnicodeCharacter{2190}{\ensuremath{\leftarrow}}
\pagestyle{plain}
\setlength{\parindent}{0pt}
\setlength{\parskip}{0pt}
\begin{document}
{\Large\bfseries Tarea de pseudocódigo}
\vspace{12pt}
\begin{enumerate}[label=\textbf{\arabic*.},leftmargin=*,itemsep=11pt,topsep=0pt]
\item \begin{minipage}[t]{\linewidth}
Según el aula, el pseudocódigo es

\vspace{4pt}
\begin{tabular}{@{}p{0.47\linewidth}p{0.47\linewidth}@{}}
a)~la descripción de un algoritmo en lenguaje natural estructurado, mezclando español con palabras clave de programación & b)~un programa ya compilado \\
c)~un lenguaje de programación que se ejecuta en la computadora & d)~un dibujo con símbolos conectados por flechas \\
\end{tabular}
\end{minipage}

\item \begin{minipage}[t]{\linewidth}
En un diagrama de flujo, ¿qué símbolo corresponde a la instrucción area $\leftarrow$ b * h?

\vspace{4pt}
\begin{tabular}{@{}p{0.47\linewidth}p{0.47\linewidth}@{}}
a)~Óvalo de inicio o fin & b)~Rectángulo de proceso \\
c)~Paralelogramo de entrada o salida & d)~Rombo de decisión \\
\end{tabular}
\end{minipage}

\item \begin{minipage}[t]{\linewidth}
En un diagrama de flujo, ¿qué forma tiene el símbolo que marca el INICIO y el FIN?

\vspace{4pt}
\begin{tabular}{@{}p{0.47\linewidth}p{0.47\linewidth}@{}}
a)~Paralelogramo & b)~Rombo \\
c)~Rectángulo & d)~Óvalo \\
\end{tabular}
\end{minipage}

\item \begin{minipage}[t]{\linewidth}
¿Qué escribe este algoritmo?

\vspace{2pt}{\small
\begin{verbatim}
INICIO
   ENTERO a, b, aux
   a ← 4
   b ← 6
   aux ← a
   a ← b
   b ← aux
   ESCRIBIR a, b
FIN
\end{verbatim}}


\vspace{4pt}
\begin{tabular}{@{}p{0.47\linewidth}p{0.47\linewidth}@{}}
a)~6, 4 & b)~4, 4 \\
c)~6, 6 & d)~4, 6 \\
\end{tabular}
\end{minipage}

\item \begin{minipage}[t]{\linewidth}
Si el usuario ingresa 15, ¿qué escribe este algoritmo?

\vspace{2pt}{\small
\begin{verbatim}
LEER edad
SI edad >= 18 ENTONCES
   ESCRIBIR "Eres mayor de edad"
FIN_SI
\end{verbatim}}


\vspace{4pt}
\begin{tabular}{@{}p{0.47\linewidth}p{0.47\linewidth}@{}}
a)~No escribe nada & b)~Eres mayor de edad \\
c)~Un mensaje de error & d)~Eres menor de edad \\
\end{tabular}
\end{minipage}

\item \begin{minipage}[t]{\linewidth}
Si el usuario ingresa 89, ¿qué escribe este algoritmo?

\vspace{2pt}{\small
\begin{verbatim}
LEER nota
SI nota >= 90 ENTONCES
   ESCRIBIR "Excelente"
SINO
   SI nota >= 70 ENTONCES
      ESCRIBIR "Aprobado"
   SINO
      ESCRIBIR "Reprobado"
   FIN_SI
FIN_SI
\end{verbatim}}


\vspace{4pt}
\begin{tabular}{@{}p{0.47\linewidth}p{0.47\linewidth}@{}}
a)~Reprobado & b)~Excelente, Aprobado \\
c)~Aprobado & d)~Excelente \\
\end{tabular}
\end{minipage}

\item \begin{minipage}[t]{\linewidth}
¿Qué escribe este algoritmo?

\vspace{2pt}{\small
\begin{verbatim}
PARA i ← 1 HASTA 3 HACER
   ESCRIBIR i * 2
FIN_PARA
\end{verbatim}}


\vspace{4pt}
\begin{tabular}{@{}p{0.47\linewidth}p{0.47\linewidth}@{}}
a)~1, 2, 3 & b)~2, 4, 6 \\
c)~2, 4, 6, 8 & d)~1, 4, 9 \\
\end{tabular}
\end{minipage}

\item \begin{minipage}[t]{\linewidth}
¿Qué escribe este algoritmo?

\vspace{2pt}{\small
\begin{verbatim}
INICIO
   ENTERO i, suma
   suma ← 0
   PARA i ← 2 HASTA 6 HACER
      suma ← suma + i
   FIN_PARA
   ESCRIBIR suma
FIN
\end{verbatim}}


\vspace{4pt}
\begin{tabular}{@{}p{0.47\linewidth}p{0.47\linewidth}@{}}
a)~6 & b)~20 \\
c)~21 & d)~18 \\
\end{tabular}
\end{minipage}

\item \begin{minipage}[t]{\linewidth}
¿Cuántas veces se ejecuta el cuerpo de este bucle?

\vspace{2pt}{\small
\begin{verbatim}
i ← 0
MIENTRAS i < 10 HACER
   i ← i + 3
FIN_MIENTRAS
\end{verbatim}}


\vspace{4pt}
\begin{tabular}{@{}p{0.47\linewidth}p{0.47\linewidth}@{}}
a)~10 & b)~5 \\
c)~3 & d)~4 \\
\end{tabular}
\end{minipage}

\item \begin{minipage}[t]{\linewidth}
¿Qué escribe este algoritmo?

\vspace{2pt}{\small
\begin{verbatim}
INICIO
   ENTERO n
   n ← 1
   MIENTRAS n < 20 HACER
      n ← n * 3
   FIN_MIENTRAS
   ESCRIBIR n
FIN
\end{verbatim}}


\vspace{4pt}
\begin{tabular}{@{}p{0.47\linewidth}p{0.47\linewidth}@{}}
a)~27 & b)~9 \\
c)~3 & d)~20 \\
\end{tabular}
\end{minipage}

\end{enumerate}
\end{document}
